% GENERATED FILE. Edit canonical lessons, then run npm run generate:books.
% canonical-source-hash: fnv1a64:f9e0bd002be8414c
% canonical-lessons: MW-C172-jharno, MW-C172-dhur, MW-C172-jalo, MW-C172-lipno, MW-C172-garo, MW-R172-review-166-172

\chapter{Sweeping and Plastering}
\label{ch:mw-sweeping-and-plastering}
\hlchaptermodality{172}

\begin{chapteropening}
\textbf{By the end of this chapter:} \emph{I can say how the house is swept and the floor plastered.}

Complete the last lesson of chapter 172: I can say how the house is swept and the floor plastered.
\end{chapteropening}

\section[\texorpdfstring{\mw{झाड़णो} (jhāṛṇo)}{झाड़णो (jhāṛṇo)}]{\mw{झाड़णो} (jhāṛṇo) — to sweep}

\label{lesson:MW-C172-jharno}



\begin{quote}
Before the new one: say the Marwadi for clean, then the Marwadi for rubbish.
\end{quote}

\subsection*{You'll want to know: \mw{झाड़णो}}
\textbf{\mw{झाड़णो}} — \emph{jhāṛṇo} — \textquotedblleft{}to sweep\textquotedblright{}.

\begin{quote}
\mw{मां} \mw{दिनूगै} \mw{कमरो} \mw{झाड़ै} \mw{है।}
\end{quote}

\begin{quote}
\emph{mā̃ dinūgai kamro jhāṛai hai.}
\end{quote}

\begin{quote}
Mother sweeps the room early in the morning.
\end{quote}

Past (I …): \mw{म्हैं} \mw{कमरो} \mw{झाड़्यो}, \emph{mhaĩ kamro jhāṛyo}.

\begin{cousinweb}
\textbf{\mw{झाड़णो}} (\emph{jhāṛṇo}) is inherited from a reconstructed \emph{*jhaṭati}, \textquotedblleft{}falls\textquotedblright{}, and its causative, \textquotedblleft{}shakes down\textquotedblright{}; CDIAL lists Marwari \emph{jhāṛṇo}, \textquotedblleft{}to sweep\textquotedblright{} (CDIAL 5328). You sweep with the broom you know, \textbf{\mw{बुहारी}}.
\end{cousinweb}

\begin{grammarlens}[title={What you've built}]
One more word for saying what you do.
\end{grammarlens}

\subsection*{Your turn}
\begin{itemize}
\raggedright
  \item \emph{Say it:} \emph{jhāṛṇo}
  \item \emph{Say it:} \emph{jhāṛṇo}, once more
  \item \emph{Say it:} \emph{mā̃ dinūgai kamro jhāṛai hai.} — Mother sweeps the room early in the morning.
  \item \emph{Recall:} say the Marwadi for clean, then the Marwadi for rubbish, then say \emph{jhāṛṇo} again
\end{itemize}

\begin{tcolorbox}[breakable,colback=teal!4,colframe=teal!35!black,title={Before you move on}]
What does \mw{झाड़णो} mean? (\textquotedblleft{}To sweep\textquotedblright{}.) Say it once more.
\end{tcolorbox}
\section[\texorpdfstring{\mw{धूड़} (dhūṛ)}{धूड़ (dhūṛ)}]{\mw{धूड़} (dhūṛ) — dust}

\label{lesson:MW-C172-dhur}



\begin{quote}
Before the new one: say the Marwadi for rubbish, then the Marwadi for to sweep.
\end{quote}

\subsection*{You'll want to know: \mw{धूड़}}
\textbf{\mw{धूड़}} — \emph{dhūṛ} — \textquotedblleft{}dust\textquotedblright{}.

\begin{quote}
\mw{मेज} \mw{माथै} \mw{धूड़} \mw{है।}
\end{quote}

\begin{quote}
\emph{mej māthai dhūṛ hai.}
\end{quote}

\begin{quote}
There is dust on the table.
\end{quote}

A feminine noun.

\begin{cousinweb}
\textbf{\mw{धूड़}} (\emph{dhūṛ}) is inherited from a reconstructed \emph{*dhūḍi}, \textquotedblleft{}dust, powder\textquotedblright{}; CDIAL lists Old Marwari \emph{dhūli} and Gujarati \emph{dhūṛ} (CDIAL 6835).
\end{cousinweb}

\begin{grammarlens}[title={What you've built}]
One more word for saying what you do.
\end{grammarlens}

\subsection*{Your turn}
\begin{itemize}
\raggedright
  \item \emph{Say it:} \emph{dhūṛ}
  \item \emph{Say it:} \emph{dhūṛ}, once more
  \item \emph{Say it:} \emph{mej māthai dhūṛ hai.} — There is dust on the table.
  \item \emph{Recall:} say the Marwadi for rubbish, then the Marwadi for to sweep, then say \emph{dhūṛ} again
\end{itemize}

\begin{tcolorbox}[breakable,colback=teal!4,colframe=teal!35!black,title={Before you move on}]
What does \mw{धूड़} mean? (\textquotedblleft{}Dust\textquotedblright{}.) Say it once more.
\end{tcolorbox}
\section[\texorpdfstring{\mw{जाळो} (jāḷo)}{जाळो (jāḷo)}]{\mw{जाळो} (jāḷo) — a cobweb}

\label{lesson:MW-C172-jalo}



\begin{quote}
Before the new one: say the Marwadi for to sweep, then the Marwadi for dust.
\end{quote}

\subsection*{You'll want to know: \mw{जाळो}}
\textbf{\mw{जाळो}} — \emph{jāḷo} — \textquotedblleft{}a cobweb\textquotedblright{}.

\begin{quote}
\mw{भींत} \mw{माथै} \mw{जाळो} \mw{है।}
\end{quote}

\begin{quote}
\emph{bhī̃t māthai jāḷo hai.}
\end{quote}

\begin{quote}
There's a cobweb on the wall.
\end{quote}

A masculine noun.

\begin{cousinweb}
\textbf{\mw{जाळो}} (\emph{jāḷo}) is inherited from Sanskrit \emph{jāla}, \textquotedblleft{}a net, a snare\textquotedblright{}; CDIAL lists Marwari \emph{jāḷ}, \textquotedblleft{}a net\textquotedblright{} (CDIAL 5213). A cobweb is the spider's net.
\end{cousinweb}

\begin{grammarlens}[title={What you've built}]
One more word for saying what you do.
\end{grammarlens}

\subsection*{Your turn}
\begin{itemize}
\raggedright
  \item \emph{Say it:} \emph{jāḷo}
  \item \emph{Say it:} \emph{jāḷo}, once more
  \item \emph{Say it:} \emph{bhī̃t māthai jāḷo hai.} — There's a cobweb on the wall.
  \item \emph{Recall:} say the Marwadi for to sweep, then the Marwadi for dust, then say \emph{jāḷo} again
\end{itemize}

\begin{tcolorbox}[breakable,colback=teal!4,colframe=teal!35!black,title={Before you move on}]
What does \mw{जाळो} mean? (\textquotedblleft{}A cobweb\textquotedblright{}.) Say it once more.
\end{tcolorbox}
\section[\texorpdfstring{\mw{लीपणो} (līpṇo)}{लीपणो (līpṇo)}]{\mw{लीपणो} (līpṇo) — to plaster (a floor or wall with mud)}

\label{lesson:MW-C172-lipno}



\begin{quote}
Before the new one: say the Marwadi for dust, then the Marwadi for a cobweb.
\end{quote}

\subsection*{You'll want to know: \mw{लीपणो}}
\textbf{\mw{लीपणो}} — \emph{līpṇo} — \textquotedblleft{}to plaster (a floor or wall with mud)\textquotedblright{}.

\begin{quote}
\mw{मां} \mw{आंगणो} \mw{लीपै} \mw{है।}
\end{quote}

\begin{quote}
\emph{mā̃ ā̃gaṇo līpai hai.}
\end{quote}

\begin{quote}
Mother plasters the courtyard.
\end{quote}

Past (I …): \mw{म्हैं} \mw{आंगणो} \mw{लीप्यो}, \emph{mhaĩ ā̃gaṇo līpyo}.

\begin{cousinweb}
\textbf{\mw{लीपणो}} (\emph{līpṇo}) is inherited from Sanskrit \emph{lipyate}, \textquotedblleft{}is smeared\textquotedblright{}; CDIAL lists Hindi \emph{līpnā}, \textquotedblleft{}to smear\textquotedblright{} (CDIAL 11061). CDIAL gives no Marwari form in this entry.
\end{cousinweb}

\begin{grammarlens}[title={What you've built}]
One more word for saying what you do.
\end{grammarlens}

\subsection*{Your turn}
\begin{itemize}
\raggedright
  \item \emph{Say it:} \emph{līpṇo}
  \item \emph{Say it:} \emph{līpṇo}, once more
  \item \emph{Say it:} \emph{mā̃ ā̃gaṇo līpai hai.} — Mother plasters the courtyard.
  \item \emph{Recall:} say the Marwadi for dust, then the Marwadi for a cobweb, then say \emph{līpṇo} again
\end{itemize}

\begin{tcolorbox}[breakable,colback=teal!4,colframe=teal!35!black,title={Before you move on}]
What does \mw{लीपणो} mean? (\textquotedblleft{}To plaster (a floor or wall with mud)\textquotedblright{}.) Say it once more.
\end{tcolorbox}
\section[\texorpdfstring{\mw{गारो} (gāro)}{गारो (gāro)}]{\mw{गारो} (gāro) — mud plaster}

\label{lesson:MW-C172-garo}



\begin{quote}
Before the new one: say the Marwadi for a cobweb, then the Marwadi for to plaster.
\end{quote}

\subsection*{You'll want to know: \mw{गारो}}
\textbf{\mw{गारो}} — \emph{gāro} — \textquotedblleft{}mud plaster\textquotedblright{}.

\begin{quote}
\mw{गारो} \mw{गीलो} \mw{है।}
\end{quote}

\begin{quote}
\emph{gāro gīlo hai.}
\end{quote}

\begin{quote}
The mud plaster is wet.
\end{quote}

A masculine noun.

\begin{cousinweb}
\textbf{\mw{गारो}} (\emph{gāro}) comes from a reconstructed \emph{*gāra}, \textquotedblleft{}dripping, oozing\textquotedblright{}; CDIAL lists Hindi \emph{gārā}, \textquotedblleft{}thick mud, plaster\textquotedblright{}, and Sindhi \emph{gāro}, \textquotedblleft{}mud plaster\textquotedblright{} (CDIAL 4137). CDIAL gives no Marwari form in this entry.
\end{cousinweb}

\begin{grammarlens}[title={What you've built}]
One more word for saying what you do.
\end{grammarlens}

\subsection*{Your turn}
\begin{itemize}
\raggedright
  \item \emph{Say it:} \emph{gāro}
  \item \emph{Say it:} \emph{gāro}, once more
  \item \emph{Say it:} \emph{gāro gīlo hai.} — The mud plaster is wet.
  \item \emph{Recall:} say the Marwadi for a cobweb, then the Marwadi for to plaster, then say \emph{gāro} again
\end{itemize}

\begin{tcolorbox}[breakable,colback=teal!4,colframe=teal!35!black,title={Before you move on}]
What does \mw{गारो} mean? (\textquotedblleft{}Mud plaster\textquotedblright{}.) Say it once more.
\end{tcolorbox}
\section[(dialogue)]{Looking back: from Up at Daybreak to Sweeping and Plastering}

\label{lesson:MW-R172-review-166-172}



\begin{quote}
Say the Marwadi for to wake up and to rinse.
\end{quote}

\subsection*{Your turn}
Say these aloud:

\begin{itemize}
\raggedright
  \item to wake up, to rinse, a night's sleep, to lie down, dusk
  \item work, first, to sprinkle, last, at once
  \item to rise, sunshine, shade, holy basil, to water plants
  \item to knead, to roll out, to boil, to fry, to taste
  \item to wash, to wring out, to spread out to dry, to dry (clothes), to gather up
  \item a utensil, a rag, to scrub, clean, rubbish
  \item to sweep, dust, a cobweb, to plaster, mud plaster
\end{itemize}

\begin{tcolorbox}[breakable,colback=teal!4,colframe=teal!35!black,title={Before you move on}]
To wake up? (\textbf{\mw{जागणो}}, \emph{jāgṇo}.) To rinse? (\textbf{\mw{खंगाळणो}}, \emph{khaṅgāḷṇo}.) A night's sleep? (\textbf{\mw{नींद}}, \emph{nī̃d}.) To lie down? (\textbf{\mw{पोढणो}}, \emph{poḍhṇo}.) Dusk? (\textbf{\mw{सांझ}}, \emph{sā̃jh}.) Work? (\textbf{\mw{काम}}, \emph{kām}.) First? (\textbf{\mw{पैलो}}, \emph{pailo}.) To sprinkle? (\textbf{\mw{छांटणो}}, \emph{chhā̃ṭṇo}.) Last? (\textbf{\mw{छेलो}}, \emph{chhelo}.) At once? (\textbf{\mw{झट}}, \emph{jhaṭ}.) To rise? (\textbf{\mw{ऊगणो}}, \emph{ūgṇo}.) Sunshine? (\textbf{\mw{धूप}}, \emph{dhūp}.) Shade? (\textbf{\mw{छांव}}, \emph{chhā̃v}.) Holy basil? (\textbf{\mw{तुळसी}}, \emph{tuḷsī}.) To water plants? (\textbf{\mw{सींचणो}}, \emph{sī̃chṇo}.) To knead? (\textbf{\mw{पोवणो}}, \emph{povṇo}.) To roll out? (\textbf{\mw{बेलणो}}, \emph{belṇo}.) To boil? (\textbf{\mw{उकळणो}}, \emph{ukaḷṇo}.) To fry? (\textbf{\mw{तळणो}}, \emph{taḷṇo}.) To taste? (\textbf{\mw{चाखणो}}, \emph{chākhṇo}.) To wash? (\textbf{\mw{धोणो}}, \emph{dhoṇo}.) To wring out? (\textbf{\mw{निचोवणो}}, \emph{nichovṇo}.) To spread out to dry? (\textbf{\mw{पसारणो}}, \emph{pasārṇo}.) To dry (clothes)? (\textbf{\mw{सुकावणो}}, \emph{sukāvṇo}.) To gather up? (\textbf{\mw{समेटणो}}, \emph{sameṭṇo}.) A utensil? (\textbf{\mw{बासण}}, \emph{bāsaṇ}.) A rag? (\textbf{\mw{चीथड़ो}}, \emph{chīthṛo}.) To scrub? (\textbf{\mw{मांजणो}}, \emph{mā̃jṇo}.) Clean? (\textbf{\mw{साफ}}, \emph{sāph}.) Rubbish? (\textbf{\mw{कचरो}}, \emph{kachro}.) To sweep? (\textbf{\mw{झाड़णो}}, \emph{jhāṛṇo}.) Dust? (\textbf{\mw{धूड़}}, \emph{dhūṛ}.) A cobweb? (\textbf{\mw{जाळो}}, \emph{jāḷo}.) To plaster? (\textbf{\mw{लीपणो}}, \emph{līpṇo}.) Mud plaster? (\textbf{\mw{गारो}}, \emph{gāro}.)
\end{tcolorbox}
